\chapter{Uncertainty, Replication, and Demarcation}
\label{ch:reliability}

\begin{chapterguide}
This chapter examines how scientific systems learn from error. It compares
Bayesian and error-statistical approaches, explains why replication is more
than repeating a number, and develops a functional account of demarcation.
The goal is not certainty, but reliable correction under uncertainty.
\end{chapterguide}

Fallibility is not a defect that science can eliminate. It is a condition that
science can manage better or worse. A reliable inquiry does not promise that
its claims are never wrong. It makes it possible to discover where they are
wrong, how much depends on them, and what should replace them.

This shift matters because scientific reliability is often discussed as if one
test could settle it. A \(p\)-value, a confidence interval, a Bayesian
posterior, a replication, or an expert consensus can each be informative. None
is self-interpreting. Objectivity lies in the relation among design, error
control, evidence, and revision.

\section{Probability as rational change}

Bayesian epistemology gives a precise rule for coherent belief revision. For
hypothesis \(H\) and evidence \(E\),

\[
  \Prob(H\mid E)
  =
  \frac{\Prob(E\mid H)\Prob(H)}{\Prob(E)}.
\]

The rule has two virtues. It forces the investigator to compare a hypothesis
with alternatives, and it distinguishes evidence from belief. The likelihood
\(\Prob(E\mid H)\) is not the same as the posterior probability
\(\Prob(H\mid E)\). A surprising result can be strong evidence for \(H\) over
\(H_0\) even when \(H\) remains improbable in an absolute sense.

The rule has limits. Priors can encode relevant background knowledge, but they
can also encode unjustified assumptions. Likelihoods depend on the model of
data collection and analysis. If an investigator tests many hypotheses and
reports only the surprising one, the reported evidence does not have the
likelihood that a preregistered test would have had.

\section{Error-statistical learning}

Deborah Mayo emphasizes learning from error: a result is informative when a
procedure would probably have produced a different result if a relevant error
were present \citep{mayo1996}. This is a frequentist and experimental
complement to Bayesian reasoning. It focuses on how a design probes error
rather than on the subjective degree of belief in a hypothesis.

The two approaches need not be enemies. Bayesian updating answers how beliefs
should change given a specified model. Error-statistical analysis asks whether
the procedure has good capability to detect errors or discrepancies. \RCR\
uses both as diagnostic tools:

\begin{itemize}
\item use likelihoods and priors to compare claims transparently;
\item use error analysis to ask what the design would have detected;
\item report model uncertainty and the possibility that the analysis was
selected after seeing the data.
\end{itemize}

\section{Replication is not a magic word}

A replication repeats a study or result under some relation of similarity.
There are several kinds:

\begin{enumerate}
\item \textbf{Exact replication:} repeats procedures as closely as possible.
\item \textbf{Conceptual replication:} tests the same claim using a different
operationalization or design.
\item \textbf{Robustness test:} changes a condition that should not matter if
the mechanism is stable.
\item \textbf{Extension:} tests whether the claim travels to a new population,
scale, or context.
\end{enumerate}

A failed replication can have several interpretations. The original effect may
be false, the replication may have lower power, the procedures may differ in a
causally relevant way, the effect may be context-sensitive, or the original
estimate may have been inflated by selection. The correct response is not to
declare a field true or false from one replication rate. It is to inspect the
conditions under which the effect appears.

The Open Science Collaboration's 2015 project attempted replications of 100
psychological studies. It reported that replicated effects were, on average,
smaller than original effects and that there was no single standard for
replication success \citep{osc2015}. The study is valuable precisely because
its interpretation requires more than one headline number. It reveals how
publication selection, statistical power, measurement, and context interact.

\section{Uncertainty and the replication ladder}

\RCR\ proposes a replication ladder. A claim receives stronger support as it
survives increasingly independent challenges:

\begin{enumerate}
\item reanalysis of the same data;
\item independent code or analytical pipeline;
\item new data with the original design;
\item different measurement or instrument;
\item different population or setting;
\item intervention on the mechanism;
\item use in a successful new prediction or application.
\end{enumerate}

The ladder is not a demand that every claim reach the final rung. A historical
claim may not permit intervention. A mathematical theorem has a different
structure. A medical treatment may be ethically tested only within a narrow
population. The ladder is a way to specify what kind of support has been
achieved, not a universal ranking of sciences.

\section{The demarcation problem}

The demarcation problem asks how to distinguish science from non-science or
pseudoscience. Popper proposed falsifiability. Lakatos emphasized progressive
problem shifts. Laudan argued that the search for a single demarcation
criterion had failed and that the more useful question is whether claims are
reliable, problem-solving, and evidentially supported \citep{laudan1977,laudan1983}.

A binary boundary is attractive for public decisions, but scientific practices
are too diverse for one property to do all the work. Some mature sciences make
few direct predictions; some non-scientific practices generate testable
predictions. A field can be methodologically serious in one period and
degenerate in another.

\RCR\ replaces a single boundary with a profile of epistemic performance. A
claim or practice is more science-like when it:

\begin{itemize}
\item creates risky contact with a target;
\item records methods, data, and uncertainty;
\item exposes itself to independent criticism;
\item updates or abandons claims when evidence requires it;
\item separates empirical results from value and metaphysical interpretation;
\item supports transportable predictions, interventions, or explanations.
\end{itemize}

This is a functional demarcation criterion. It does not say that science is
whatever produces useful outcomes. It says that the institutional practice
deserves scientific authority to the extent that it maintains correction
capacity under world-sensitive constraints.

\section{Scientism and the misuse of science}

Scientism is not simply confidence in science. It is the unwarranted extension
of scientific authority beyond the question, method, or evidence that supports
it. A physical law does not by itself determine an ethical obligation. A
statistical regularity does not settle a political priority. A model of human
behaviour does not eliminate interpretation, meaning, or agency.

\RCR\ rejects scientism for the same reason it rejects relativism: both
collapse distinctions. Science is authoritative for claims that its methods
can constrain. It is not a universal replacement for history, ethics,
political judgment, or lived interpretation. This limitation protects science
from being asked to prove what its methods cannot address, and protects other
forms of reasoning from being dismissed by an illegitimate appeal to scientific
status.

\section{Reliability as a system property}

A single true claim can emerge from an unreliable process by luck. A single
false claim can emerge from a generally reliable process by chance. The
objectivity of science is therefore not exhausted by the truth value of
individual outputs. It includes the design of a system that tends to correct
errors over time.

A system's reliability depends on:

\begin{enumerate}
\item the quality of its measurements and models;
\item the independence and competence of its critics;
\item the incentives attached to novelty and correction;
\item the visibility of negative results and failed replications;
\item the capacity to revise concepts, not only parameter estimates.
\end{enumerate}

The fifth point is crucial. A field can replicate a measurement reliably while
measuring the wrong construct. Objectivity sometimes requires changing the
question.

\begin{principlebox}
\textbf{Correction principle.} The epistemic standing of a scientific claim
depends not only on present support but on the quality of the process that
could reveal its failure. A claim with moderate evidence and strong correction
capacity can be more trustworthy than a claim with impressive evidence inside
an opaque and uncorrectable system.
\end{principlebox}

\section*{Chapter summary}

Probability formalizes belief revision; error statistics asks what a procedure
could have detected; replication tests stability under controlled variation;
and demarcation is best understood as a profile of correction capacity rather
than one binary test. Scientific authority should be proportional to the
quality of world-sensitive error correction. This supports fallibilism without
relativism and confidence without scientism.

\begin{discussion}
\begin{enumerate}
\item When does a failed replication count against a claim?
\item Is a transparent but imperfect method preferable to a highly accurate
method whose assumptions cannot be inspected?
\item Can a practice be scientific in one domain and non-scientific in another?
\end{enumerate}
\end{discussion}

\begin{furtherreading}
See Mayo (1996), the Open Science Collaboration (2015), Ioannidis (2005),
Laudan (1977, 1983), and Reiss and Sprenger (2014).
\end{furtherreading}

